\chapter{Formal Methods with Lean 4}

A file named \texttt{Sertifikat.lean} arrived together with a single sentence of introduction:
``everything has already been proved by machine''. It contained forty one-line theorems, almost all
of them ending with \texttt{by native\_decide}, and at the bottom there was one summary constant
that mentioned the threshold $10^{-3}$. Not one of its lines can we inspect with our eyes in a
reasonable amount of time: what is being compared are sixteen-digit integers that come from
turning binary fractions into exact form. That introductory sentence is true in a narrow sense and
misleading in an important sense. The machine does check every line, but it does not check whether
what is being checked is what was meant.

Testing and proving promise different things, and the difference is not a question of degree of
confidence. Testing takes a number of inputs, runs the program, and then reports what happened on
those inputs. The conclusion always has the form ``for $10^{4}$ test samples, nothing failed''. A
formal proof promises more: one statement quantified over all inputs, together with an object that
a third party can check without trusting its author. That promise is kept only when three
conditions hold at once, namely that the statement is written correctly, that the object being
checked really is a proof of that statement, and that the foundation of the checking has not been
misused. This chapter takes those three conditions apart on one concrete system, Lean 4.

The tools of the previous chapter are not enough for this work. Chapter 27 taught us how to report
what is actually measured: keep the training data apart from the test data, state the distribution,
report the variance across seeds. That is a discipline about \emph{empirical claims}. Here we face a
claim of a different shape: not an average over samples, but a statement quantified as
$\forall x,\ P(x)$. For a claim of that kind no confidence interval helps, because a single
counterexample is enough to bring it down, and the absence of a counterexample is not a proof. What
is needed is an object that can be rechecked by someone who does not believe us.

\begin{untukanak}
Imagine a very meticulous chess referee. He never asks whether the game is interesting, whether
both players play well, or whether the rules of chess are in fact the rules that ought to be used.
His task is only one: every time a move appears, he checks whether that move is legal according to
the rule book. If all the moves are legal, he signs the board and says ``this game is valid''. He
does not care who wins.

Now imagine someone hands over a board full of legal moves, then says ``this proves that I can beat
anyone''. The referee will answer honestly: what I checked is only that the moves are legal, not
that your sentence is true. If the sentence is wrong, a legal board does not help at all, because
what is wrong is not the moves but the question.

Have you ever given a correct answer to the wrong problem? Or done your homework neatly when what
was asked for was a different page? In this chapter we will find out why a computer can check a
proof down to the smallest step, why that checking does not in the least guarantee that the
question is right, and how people write the right question in a language that a machine can check.
\end{untukanak}

\section{Informal proof and formal proof}

\tingkat{intermediate}

A proof in a mathematics book is a text that convinces a trained reader. It contains jumps that are
agreed on in common: ``clearly'', ``without loss of generality'', ``analogous for the other cases''.
Those jumps are not cheating; they are a compression accepted by the community. The trouble is that
such compression cannot be checked by a machine, and cannot be separated from the knowledge of its
reader.

A formal proof is another version of the same object. Its statement is written in a language with a
strict grammar and strict types, and its proof is a \emph{term}, namely an expression whose type is
that statement. The checking is carried out by the kernel of the system, a small piece of code that
knows only a few rules and knows nothing at all about mathematics. Because the kernel is small, its
checking can be trusted far more strongly than any large program around the system.

\kamus{formal verification}{verifikasi formal}{Step-by-step checking by a clerk who is never bored
and never trusts the author.}

\begin{definisi}[Formal proof]
Let $P$ be a statement written in the formal language of the system. A \emph{formal proof} of $P$ is
a term $p$ such that $p:P$ can be derived by the kernel of the system from definitions, theorems
that have already been checked, and axioms declared open. The system accepts $P$ if and only if a
term of that kind is found and passes the type check.\\index{formal verification}
\end{definisi}

That definition contains three words that are easy to pass over. ``Can be derived'' means that
there is a mechanical procedure, not that there is a person who agrees. ``Checked'' means that the
outcome is binary: accepted or rejected, with no degrees of confidence. ``Declared open'' means that
the axiom is recorded, not hidden; we will use that recording later as our main instrument.

The difference that is most often misunderstood is the difference between \emph{truth} and
\emph{provability}. Lean does not know whether a statement is true in the real world. It knows only
whether there is a term whose type is that statement. If the statement is false, that term will
never be found, but the absence of a term is not a proof that the statement is false; it may simply
be that we have not found it yet.

\begin{bacakeras}
First voice: ``I have checked it, everything is correct.''
Second voice: ``Show me how.''
First voice: ``I read it again three times.''
Second voice: ``That is a habit, not a method.''
First voice: ``I am sure.''
Second voice: ``Certainty is not proof. Give me the steps, one by one.''

That is what a conversation with a machine is like. It does not ask whether we are good people,
whether we have had breakfast, or whether we graduated with high marks. It asks one thing,
patiently, every time, with no angry tone: show me the steps. If the steps are there, the door
opens. If the steps are not there, the door stays shut, as tightly as before, even if we repeat the
request a hundred times. The machine does not judge our intentions; it only checks the form.

Say it together, slowly, then repeat it once more: the machine does not believe, it checks.
\end{bacakeras}

\section{Propositions as types, proofs as terms}

\tingkat{intermediate}

In Lean there is no gulf between a statement and its proof. Both are terms, checked by the same
rules. Every sentence that can take the value true or false is given the type $\mathtt{Prop}$; a
proof of that sentence is a term whose type is that sentence. The statement ``$P$ is proved''
therefore does not mean that $P$ is true in some model, but that there is a term of type $P$ that
passes the check.

\kamus{proposition}{proposisi}{A sentence that can be judged true or false, like an exam question
whose answer must be proved, not merely chosen.}

\begin{definisi}[Inhabitant of a type]
A type $T$ is called \emph{inhabited} when there is a term $t$ with $t:T$; the notation $t:T$ is
read ``$t$ has type $T$''. If $P:\mathtt{Prop}$, then $P$ is called a \emph{proposition} and $t:P$
is called a \emph{proof term} for $P$.\\index{proof term} Two terms are called \emph{definitionally
equal} when both reduce to an identical form by symbolic computation.\\index{proposition-type}
\end{definisi}

This dictionary, namely the correspondence between logic and types, is known as
Curry--Howard.\\index{Curry-Howard} Implication becomes a function type, because a proof of
$P\Rightarrow Q$ must turn an arbitrary proof of $P$ into a proof of $Q$, and that is the work of a
function. The universal quantifier becomes a dependent function, because a proof of
$\forall x,\ P(x)$ must produce a proof of $P(a)$ for every $a$ that is given. Conjunction becomes
a pair, because a proof of $P\wedge Q$ contains both of its proofs at once. The last row of the
table below is the most important: $\bot$ has no inhabitant.

\begin{center}
\begin{tabular}{lll}
\hline
Logic & Type in Lean & How to build the term \\
\hline
$P\Rightarrow Q$   & $P\to Q$              & \texttt{fun h => ...} \\
$P\wedge Q$        & \texttt{And P Q}      & \texttt{And.intro} \\
$P\vee Q$          & \texttt{Or P Q}       & \texttt{Or.inl}, \texttt{Or.inr} \\
$\top$             & \texttt{True}         & \texttt{True.intro} \\
$\bot$             & \texttt{False}        & no inhabitant \\
$\forall x,\ P(x)$ & $\forall x,\ P\,x$    & \texttt{fun x => ...} \\
$\exists x,\ P(x)$ & \texttt{Exists}       & \texttt{Exists.intro} \\
\hline
\end{tabular}
\end{center}

\kamus{proof term}{term bukti}{A building frame whose shape follows the blueprint of its statement
exactly; if the shape does not fit, the building is rejected.}

The following terms are the smallest useful proofs: the identity function as a proof of the identity
implication, and the composition of two functions as a proof of the transitivity of implication.
Both are checked without using any knowledge at all about the contents of $P$, $Q$, and $R$.

\begin{verbatim}
-- proof of the identity implication, term style
def bukti_id {P : Prop} : P -> P := fun h => h

-- composition of two implications
def bukti_komposisi {P Q R : Prop} (h1 : P -> Q) (h2 : Q -> R) : P -> R :=
  fun hp => h2 (h1 hp)

#check bukti_id
#check bukti_komposisi
#check Prop
\end{verbatim}

The output of the last command, as it stands, is \texttt{Prop : Type}: the set of propositions is
itself a type. The two lines before it print the types of the two definitions. Note that
\texttt{\#check} does not prove anything; it only asks for the type, and that is also what the
kernel of the system does when it checks.

\begin{proposisi}
For every $P:\mathtt{Prop}$ we have $P\to P$. For $P,Q,R:\mathtt{Prop}$ we have
$(P\to Q)\to(Q\to R)\to P\to R$.
\end{proposisi}

\begin{proof}
The term $\lambda h.\,h$ has type $P\to P$: in the context $h:P$, the variable $h$ itself is a
proof of $P$, so abstraction over $h$ produces a function from $P$ to $P$. For the second statement,
let $h_{1}:P\to Q$, $h_{2}:Q\to R$, and $h_{p}:P$. The application $h_{1}\,h_{p}$ has type $Q$, then
the application $h_{2}\,(h_{1}\,h_{p})$ has type $R$. Three abstractions in succession over $h_{p}$,
$h_{2}$, and $h_{1}$ give a term of type $(P\to Q)\to(Q\to R)\to P\to R$.
\end{proof}

The tree in Figure~\ref{fig:pohon-tipe} writes that second proof as a derivation tree: three leaves,
two application stems, one root. The rightmost leaf is a hypothesis that is given, not a result of
derivation.

\begin{figure}[htbp]
\centering
\begin{tikzpicture}[font=\small,
  daun/.style={draw=layergrey,fill=layergrey!6,rounded corners=2pt,inner sep=4pt},
  hip/.style={draw=layergrey,dashed,fill=layergrey!3,rounded corners=2pt,inner sep=4pt},
  batang/.style={thick,layergrey}]
  \node[daun] (h1) at (-3.7,2.7) {\texttt{h1 : P -> Q}};
  \node[daun] (hp) at (-0.6,2.7) {\texttt{hp : P}};
  \node[daun] (ap1) at (-2.2,1.35) {\texttt{h1 hp : Q}};
  \node[hip]  (h2) at (2.7,1.35) {\texttt{h2 : Q -> R}};
  \node[daun] (akar) at (0,0) {\texttt{h2 (h1 hp) : R}};
  \draw[batang] (h1.south) -- (ap1.north west);
  \draw[batang] (hp.south) -- (ap1.north east);
  \draw[batang] (ap1.south) -- (akar.north west);
  \draw[batang] (h2.south) -- (akar.north east);
\end{tikzpicture}
\caption{The type tree for the term \texttt{h2 (h1 hp)}: three leaves and two application stems, with
a root of type $R$. The dashed boxes are hypotheses that are given, not results of derivation.}
\label{fig:pohon-tipe}
\end{figure}

There is one property that must be mentioned because it explains why proofs are used as
certificates rather than as data. The Lean kernel treats any two proofs of the same proposition as
definitionally equal. This property is called \emph{proof irrelevance}, and its consequence is that
a proof carries no information other than ``this statement is true''. For that reason a proof can be
thrown away once it has been checked, and what remains is only the status of the statement.

One warning is installed from the start. The Curry--Howard correspondence holds in full for
constructive logic. As soon as we use the law of excluded middle, $\mathtt{Classical.em}:P\vee\neg
P$, we add an axiom, and a proof that uses it stops being computational: it remains valid, but it is
no longer read as a program that produces a witness. Many theorems of the large libraries carry the
axiom of choice in their list of axioms, and the last part of this chapter shows how to ask Lean for
that list directly.

\section{Quantifiers: introduction and elimination rules}

\tingkat{intermediate}

An interesting statement is almost always quantified.\\index{quantifier} ``For every natural
number'' and ``there is a natural number'' are not decorations of language; both change the kind of
object that must be supplied. In the type view, that difference becomes the difference between a
function and a pair.

\kamus{quantifier}{kuantor}{The words ``every'' and ``there is''; the first demands a recipe for all
cases, the second demands one example together with its proof.}

The rule for the universal quantifier is entirely the same as the rule for a function. To prove
$\forall x,\ P(x)$, we take an arbitrary $x$ and prove $P(x)$ for that $x$; that is abstraction. To
use the hypothesis $\forall x,\ P(x)$, we apply that hypothesis to an object; that is application.
The rule for the existential quantifier is the rule for a pair: to prove $\exists x,\ P(x)$ we
supply one witness $a$ together with a proof of $P(a)$; to use an existential hypothesis we unpack
that pair into the witness and its proof.

\begin{center}
\begin{tabular}{lll}
\hline
Form of the statement & How to prove it & How to use it \\
\hline
$\forall x,\ P(x)$      & \texttt{intro x}              & apply to an object \\
$\exists x,\ P(x)$      & \texttt{exact Exists.intro a h} & \texttt{cases} on the hypothesis \\
$P\wedge Q$             & \texttt{constructor}          & \texttt{h.1}, \texttt{h.2} \\
$P\vee Q$               & \texttt{left}, \texttt{right} & \texttt{cases} on the hypothesis \\
$\neg P$                & \texttt{intro h} then reach $\bot$ & apply to a proof of $P$ \\
\hline
\end{tabular}
\end{center}

\begin{tangan}
Prepare ten blank cards and one marker. On five cards write introduction rules: ``take an
arbitrary'', ``supply a witness'', ``pair two proofs'', ``choose left'', ``choose right''. On the
other five cards write use rules: ``apply to an object'', ``unpack the pair'', ``take the first
component'', ``split into cases'', ``use it on the opponent's proof''.

Lay one goal card at the end of the table, for example ``every natural number is zero or a
successor''. Build a chain toward the goal card: the first card must be an introduction rule whose
shape fits the goal, not a use rule. Every following card may only work on the part that is still
open. A card that has no remainder to work on is pushed to the edge of the table; in Lean that card
will never be used.

What you must see: the chain gets shorter toward the goal, and there are never two cards working on
the same remainder. If that happens, your chain branches, and the machine demands two proofs, not
one.
\end{tangan}

The following example shows both at work on the natural numbers. The first statement is a tautology
that requires no reasoning at all; the second statement says that every natural number is zero or
the successor of something, and its proof requires case splitting on the shape of the number.

\begin{verbatim}
-- universal quantifier: take an arbitrary n, then compute
theorem nol_kanan : forall n : Nat, n + 0 = n := by
  intro n
  rfl

-- existential quantifier: supply a witness, then prove its property
theorem ada_tiga : exists n : Nat, n = 3 := by
  exact Exists.intro 3 rfl

-- every natural number is zero or a successor
theorem nol_atau_suksesor (n : Nat) : n = 0 \/ exists m : Nat, n = m + 1 := by
  cases n with
  | zero => left; rfl
  | succ k => right; exact Exists.intro k rfl
\end{verbatim}

The three theorems are checked without any additional axiom, and that can be confirmed with the
command \texttt{\#print axioms}, which we will use many times. Note also that
\texttt{Exists.intro k rfl} uses \texttt{rfl} as a proof of equality: in the case
\texttt{succ k}, the goal is $k+1=k+1$, which is true by definition.

\begin{catatan}
The reverse direction does not hold automatically. From $\exists x,\ P(x)$ we may not conclude
$\forall x,\ P(x)$, and Lean will reject every attempt in that direction. The mistake of swapping
the order of quantifiers is the most frequent specification error of all, and it slips through
testing without complaint because all test samples satisfy it. The part on limits at the end of this
chapter discusses its consequences.
\end{catatan}

\section{Tactics: turning the goal into a term}

\tingkat{basic}

Writing proof terms directly quickly becomes unreadable, because a deep term has to be read from the
inside out. That is why Lean provides a tactic mode.\\index{tactic} After the word \texttt{by}, we
do not write a term, but a sequence of commands that change the \emph{goal state}. The goal state is
printed as a list of hypotheses, then a separator line, then the target.

\kamus{tactic}{taktik}{A step-by-step cooking instruction, not a complete recipe that has to be
memorized all at once.}

\begin{verbatim}
-- the goal state is printed after every tactic; the line "|-" reads "therefore"
theorem imp_trans {P Q R : Prop} (h1 : P -> Q) (h2 : Q -> R) : P -> R := by
  intro hp
  trace_state          -- P Q R : Prop, h1 : P -> Q, h2 : Q -> R, hp : P, |- R
  apply h2
  trace_state          -- the target becomes |- Q
  exact h1 hp
\end{verbatim}

The output of \texttt{trace\_state} on this test machine is the following two blocks. Lean prints
the arrow as a Unicode character; here everything is rewritten in ASCII so as to be consistent with
the whole chapter.

\begin{verbatim}
P Q R : Prop
h1 : P -> Q
h2 : Q -> R
hp : P
|- R

P Q R : Prop
h1 : P -> Q
h2 : Q -> R
hp : P
|- Q
\end{verbatim}

The tactic \texttt{intro} takes a target of implication shape and moves its antecedent into the
hypotheses; that is exactly function abstraction. The tactic \texttt{exact} closes the target with a
term whose type fits exactly. The tactic \texttt{apply} matches the conclusion of a lemma against
the target, then leaves the remaining arguments as new goals. These three tactics are already enough
for all the pure-logic proofs in this chapter, and all three are surface forms of the same rule.

Figure~\ref{fig:dua-keadaan} compresses those two states into two boxes. What changes is only the
line below the bar: the target shrinks from $R$ to $Q$. The list of hypotheses does not grow by a
single line, because \texttt{apply} does not create a new proof; it only lowers the demand.

\begin{figure}[htbp]
\centering
\begin{tikzpicture}[font=\small,
  kotak/.style={draw=layergrey,fill=layergrey!6,rounded corners=3pt,align=left,inner sep=6pt,
                text width=3.9cm},
  panah/.style={-{Stealth[length=2.2mm]},thick,layerteal}]
  \node[kotak] (awal) at (0,0) {\textbf{state 1}\\[3pt]
    \texttt{P Q R : Prop}\\
    \texttt{h1 : P -> Q}\\
    \texttt{h2 : Q -> R}\\
    \texttt{hp : P}\\[3pt]
    \rule{3.6cm}{0.4pt}\\
    $\vdash$ \texttt{R}};
  \node[kotak] (akhir) at (6.7,0) {\textbf{state 2}\\[3pt]
    \texttt{P Q R : Prop}\\
    \texttt{h1 : P -> Q}\\
    \texttt{h2 : Q -> R}\\
    \texttt{hp : P}\\[3pt]
    \rule{3.6cm}{0.4pt}\\
    $\vdash$ \texttt{Q}};
  \draw[panah] (awal.east) -- node[above,font=\scriptsize]{\texttt{apply h2}} (akhir.west);
\end{tikzpicture}
\caption{The two goal states in the proof of \texttt{imp\_trans}: the hypotheses stay at four lines,
the target shrinks from $R$ to $Q$ after one tactic. The second box is still open; it is closed by
\texttt{exact h1 hp}.}
\label{fig:dua-keadaan}
\end{figure}

\begin{center}
\begin{tabular}{lll}
\hline
Tactic & What it does & When to use it \\
\hline
\texttt{intro}     & antecedent becomes a hypothesis   & target of shape $P\to Q$ \\
\texttt{exact}     & closes the target with a term    & the term is already available \\
\texttt{apply}     & lemma conclusion matches target & the lemma ends the same way \\
\texttt{rw}        & rewrites with an equation & there is a hypothesis $a=b$ \\
\texttt{simp}      & simplifies repeatedly      & many trivial steps \\
\texttt{cases}     & splits the target by constructor & case analysis is needed \\
\texttt{induction} & base case and inductive step     & recursive target over \texttt{Nat} \\
\hline
\end{tabular}
\end{center}

The tactic \texttt{rw} rewrites the target using an equation. It works in the direction that is
given, and that direction matters: rewriting in the wrong direction will not find the pattern that
is being looked for. Note the two attempts below, one of which succeeds and one of which is
rejected.

\begin{verbatim}
theorem geser (a b : Nat) (h : a = b) : a + 1 = b + 1 := by
  rw [h]        -- the target becomes b + 1 = b + 1, then it is closed by definition
  rfl

-- commutativity of addition does NOT hold by definition, so rfl fails
theorem tukar_salah (n : Nat) : n + 1 = 1 + n := by
  rfl           -- error: The rfl tactic failed

-- the correct version: rewrite with the commutativity lemma
theorem tukar_benar (n : Nat) : n + 1 = 1 + n := by
  rw [Nat.add_comm]
\end{verbatim}

The failure in the middle is a lesson, not an accident. Addition on the natural numbers is defined
by recursion on the second argument, so $n+1$ reduces to $\mathtt{succ}\,n$ whereas $1+n$ reduces to
$\mathtt{succ}\,(0+n)$, and the two stop at different forms. What we call the ``commutative
property'' at school is a theorem, not a definition, and Lean refuses to treat it as if it were a
definition. This is the most noticeable difference between reading and writing formal proofs: every
thing that we usually skip over has to be named.

The tactic \texttt{cases} splits the target according to the constructors of the type. On the
natural numbers it gives two cases, zero and successor. The tactic \texttt{simp} closes a goal by
rewriting with a collection of lemmas that are marked; it is useful but also dangerous, because its
output depends on the collection of lemmas that is currently active. A proof that begins with
\texttt{simp} without arguments can stop working when the library changes, and that is why in this
chapter \texttt{simp} is always mentioned together with the list of lemmas it uses.

One warning about tactics that are not here. Names such as \texttt{linarith}, \texttt{ring},
\texttt{norm\_num}, and \texttt{induction'} come from the Mathlib library, not from the Lean kernel.
All the code fragments in this chapter were run with Lean 4.10.0 without Mathlib, so those tactics
are not available and are not used; in their place we use \texttt{omega} (which is in the kernel)
and \texttt{simp} with an explicit list of lemmas. This difference is not a matter of style: the
library adds thousands of lemmas and several tactics, and every addition also determines what can be
closed in one line.

\begin{untukanak}
Take a pile of small stones. You arrange one stone into the smallest square. Then you add three
stones around it, and the result is a larger square. Add five more stones, and the result is a
larger square still. Add seven, nine, eleven, and so on, always increasing by two, and every time
the result is still a tidy square.

Now the question: will this go on being so if we add up to a hundred times, up to a million times?
You cannot try all of them. Nor do you need to. If you can show that the pattern of adding stones
always closes one corner and one side of the square neatly, then one small explanation holds for
every size, including sizes that you will never be able to arrange by hand.

The lazy computer does not want to try a million cases, and the overconfident computer does not want
to check anything. In this chapter we will find out why there is a third way: explain the pattern
once, then ask the machine to check that explanation down to its last letter.
\end{untukanak}

\section{A complete inductive proof: the sum of odd numbers}

\tingkat{advanced}

We are going to prove one arithmetic statement that can be checked entirely by machine, namely that
the sum of the first $n$ odd numbers equals $n^{2}$. This statement was chosen for three reasons:
the formula can be computed with Python without Lean, the proof requires induction, and the
algebraic step is small enough that every tactic can be accounted for.

\kamus{induction}{induksi}{Explaining a pattern with only two things: the smallest case, and how to
climb one step from any case.}

Before writing code, we state the statement and its proof in ordinary language. The function is
defined recursively: the sum of the first zero odd numbers is zero, and the sum of the first $n+1$
odd numbers is the sum of the first $n$ odd numbers plus the odd number number $n+1$, namely
$2n+1$.

\begin{definisi}[Sum of odd numbers]
The function $S:\mathbb{N}\to\mathbb{N}$ is defined by $S(0)=0$ and $S(n+1)=S(n)+(2n+1)$ for every
$n\in\mathbb{N}$.\\index{induction}
\end{definisi}

The number $2n+1$ is not an arbitrary choice: the $n$th odd number really does have that form, and
$2n+1$ is also the difference between $(n+1)^{2}$ and $n^{2}$. That last equality is the heart of the
proof.

\begin{lemma}
For every $k\in\mathbb{N}$ we have $(k+1)^{2}=k^{2}+2k+1$.
\end{lemma}

\begin{proof}
$(k+1)^{2}=(k+1)(k+1)=k(k+1)+(k+1)=k^{2}+k+k+1=k^{2}+2k+1$. Each step uses distributivity and
associativity of addition on the natural numbers.
\end{proof}

\begin{teorema}
For every $n\in\mathbb{N}$ we have $S(n)=n^{2}$.
\end{teorema}

\begin{proof}
By induction on $n$. For $n=0$, both sides have the value $0$. Suppose the statement is true for
$k$, namely $S(k)=k^{2}$. Then $S(k+1)=S(k)+(2k+1)=k^{2}+2k+1=(k+1)^{2}$, with the lemma above at
the last step. So the statement is true for $k+1$. By the principle of induction, the statement is
true for every $n\in\mathbb{N}$.
\end{proof}

The proof above is a complete informal proof, in the sense that no step is skipped. Now we write it
for the machine. Note that the recursive definition and the theorem are almost word for word what we
have just written, whereas the proof turns into a sequence of tactics.

\begin{verbatim}
-- recursive definition: two equations, without ":=" on the first line
def oddSum : Nat -> Nat
  | 0     => 0
  | n + 1 => oddSum n + (2 * n + 1)

theorem oddSum_sq (n : Nat) : oddSum n = n * n := by
  induction n with
  | zero => rfl
  | succ k ih =>
    rw [oddSum]        -- unfold the definition: oddSum (k+1) = oddSum k + (2*k+1)
    rw [ih]            -- use the induction hypothesis: oddSum k = k*k
    simp [Nat.succ_eq_add_one, Nat.mul_succ, Nat.add_mul, Nat.one_mul]
    omega              -- close the linear remainder
\end{verbatim}

The line \texttt{induction n with} generates two goals. The first goal is labelled \texttt{zero} and
reads $S(0)=0\cdot 0$, and it is closed by \texttt{rfl} because both sides reduce to zero. The
second goal is labelled \texttt{succ k} and carries one extra hypothesis named \texttt{ih}, namely
$S(k)=k\cdot k$. The name \texttt{ih} is chosen by us; its role is what matters.

The line \texttt{rw [oddSum]} unfolds the definition on the left-hand side. After that line the goal
reads $S(k)+(2k+1)=(k+1)(k+1)$. The line \texttt{rw [ih]} replaces $S(k)$ with $k\cdot k$, so the
goal becomes $k\cdot k+(2k+1)=(k+1)(k+1)$.

At this point we stand exactly where the lemma above does its work, but the machine does not yet know
that $(k+1)(k+1)$ can be written as $k\cdot k+2k+1$. \texttt{simp} with four explicit lemmas changes
the shape of the right-hand side into $k\cdot k+k+(k+1)$, which is a shape that is entirely linear in
$k$. The last line, \texttt{omega}, closes a linear goal of that kind.

It must be stressed that \texttt{omega} is not magic and is not a hidden axiom either. It is a
decision procedure for linear arithmetic over the integers, and it accepts only goals all of whose
operations are linear. If we call it before \texttt{simp}, the goal still contains both $k\cdot k$
and $(k+1)(k+1)$. On this test machine, the temporary experiment file named \texttt{odd.lean}
produces a rejection:

\begin{verbatim}
odd.lean:11:4: error: omega could not prove the goal:
a possible counterexample may satisfy the constraints
  a - b + 2*c <= -2
  ...
\end{verbatim}

That message is honest and useful: it names a possible counterexample, and precisely because of that
we know that what fails is the shape of the goal, not the truth of the statement. The order of
tactics is therefore not a matter of taste; it changes the kind of goal that the decision procedure
faces.

\begin{catatan}
If \texttt{simp} above is run without \texttt{omega}, the remaining goal has the shape
$k\cdot k+(1+2k)=k+(k+(k\cdot k+1))$. Mathematically that is true, but it is no longer linear, and
the simplifier does not close it. On this test machine that experiment ends with
\texttt{unsolved goals}. This is the most common shape of a failed formal proof: not that the
statement is false, but that the shape of the goal does not fit the tool in use.
\end{catatan}

\begin{akibat}
For every $n\in\mathbb{N}$, the sum of the first $n$ odd numbers is a perfect square.
\end{akibat}

That corollary is only a repetition of the theorem in other words, and that is precisely its use: it
restates the result in the language used in this part, so that a reader who meets the formula
elsewhere can recognize it.

\section{Structures, typeclasses, and what is really guaranteed}

\tingkat{advanced}

A certificate is rarely a single loose theorem. Usually it is a collection of data together with
claims about that data: thresholds, tables, file names, and several statements. Lean provides two
tools for grouping things of this kind, namely \emph{structure} and \emph{typeclass}.

\kamus{typeclass}{type class}{A list of capabilities that a type must have in order to be usable in
a certain place, like the licensing requirement for driving a particular kind of vehicle.}

\begin{definisi}[Structure and typeclass]
A \emph{structure} is a type whose fields are named, so that a value of it is built by naming the
fields. A \emph{typeclass} is a structure that is used to look for an \emph{instance}
automatically: when a definition requires a certain capability from a type, Lean searches for a
matching instance in the environment that is currently active.\\index{typeclass}
\end{definisi}

The practical difference is simple. A structure is used when the data itself is what matters, for
example a certificate with a claim field and a threshold field. A typeclass is used when what
matters is the capability: something that can be used as a proof, something whose norm can be
computed, something that can be compared. In both, the checking remains equally strict: every field
must be filled, and every instance must pass the type check.

\begin{verbatim}
-- structure: data with named fields
structure Sertifikat where
  klaim  : String
  ambang : Nat

-- typeclass: a capability that is looked up automatically
class Terperiksa (a : Type) where
  jumlah_teorema : a -> Nat

instance : Terperiksa Sertifikat where
  jumlah_teorema _ := 1

def contoh : Sertifikat := { klaim := "oddSum", ambang := 100 }

#eval (Terperiksa.jumlah_teorema contoh)   -- prints 1
\end{verbatim}

Now the part that is most often skipped in popular talk about Lean, namely the question of what the
acceptance of a theorem actually guarantees. Four things are guaranteed, and four other things are
not.

Figure~\ref{fig:alur-periksa} puts those three stages on one line: the statement is written, the
term is submitted, the kernel decides. There are only two outputs; there is no third box that says
``perhaps''.

\begin{figure}[htbp]
\centering
\begin{tikzpicture}[font=\small,
  tahap/.style={draw=layergrey,fill=layergrey!6,rounded corners=3pt,align=center,inner sep=5pt,
                text width=2.9cm,minimum height=1.7cm},
  panah/.style={-{Stealth[length=2.2mm]},thick,layergrey},
  terima/.style={draw=layerjade,fill=layerjade!8,rounded corners=3pt,align=center,inner sep=5pt,
                 text width=2.9cm},
  tolak/.style={draw=accent,fill=accent!6,rounded corners=3pt,align=center,inner sep=5pt,
                text width=2.9cm}]
  \node[tahap] (s1) at (0,0) {\textbf{1. statement}\\$P$ has type \texttt{Prop}};
  \node[tahap] (s2) at (3.9,0) {\textbf{2. term}\\$p$ is submitted for $P$};
  \node[tahap] (s3) at (7.8,0) {\textbf{3. kernel}\\check whether $p:P$};
  \draw[panah] (s1) -- (s2);
  \draw[panah] (s2) -- (s3);
  \node[terima] (ya) at (5.4,-2.9) {\textbf{accepted}\\there is a term of type $P$};
  \node[tolak] (tidak) at (9.9,-2.9) {\textbf{rejected}\\there is no term};
  \draw[panah] (s3.south west) -- node[above left,font=\scriptsize]{matches} (ya.north);
  \draw[panah] (s3.south east) -- node[above right,font=\scriptsize]{does not match} (tidak.north);
\end{tikzpicture}
\caption{The checking flow: three stages (statement, term, kernel) and two outputs, accepted or
rejected. There is no third output; a file that stops in the middle is reported as a goal that is
not yet closed, not as ``almost correct''.}
\label{fig:alur-periksa}
\end{figure}

What is guaranteed: (i) there is a term whose type is exactly the statement that was written; (ii)
that term is checked by the kernel, not by the tactic that produced it; (iii) the checking uses only
definitions, other theorems that have already been checked, and recorded axioms; (iv) the list of
axioms that are used can be asked for at any time with \texttt{\#print axioms}. Point (iv) is the
most useful instrument in this whole chapter, because it turns ``just believe me'' into ``check it
yourself''.

What is not guaranteed: (i) that the statement says what we mean; (ii) that the definitions model
objects in the real world; (iii) that the kernel, the parser, and the hardware are free of errors;
(iv) that the tactics in use do not add hidden assumptions. Point (iv) is not a theoretical flaw: it
can be checked, and the result is often surprising.

\begin{contoh}[Two ways to prove the same thing, two lists of axioms]
Consider two proofs of the statement $S(10)=100$, both of them correct. The first uses
\texttt{decide}, which runs the computation inside the kernel. The second uses
\texttt{native\_decide}\\index{native-decide}, which compiles the computation into machine code and
runs it. The two statements are the same, but their lists of axioms differ, and that difference is
visible from the command \texttt{\#print axioms}.
\end{contoh}

\begin{verbatim}
def oddSum : Nat -> Nat
  | 0     => 0
  | n + 1 => oddSum n + (2 * n + 1)

theorem cek_native : oddSum 10 = 100 := by native_decide
theorem cek_kernel : oddSum 10 = 100 := by decide

#print axioms cek_native
#print axioms cek_kernel
\end{verbatim}

The output on this test machine, as it stands:

\begin{verbatim}
'cek_native' depends on axioms: [Lean.ofReduceBool]
'cek_kernel' does not depend on any axioms
\end{verbatim}

The first line is the heart of the whole discussion about \texttt{native\_decide}. The computation
is not carried out by the kernel, but by compiled code, and Lean records that trust as an axiom
named \texttt{Lean.ofReduceBool}. In practice this is not a disaster: the Lean compiler is used to
compile Lean itself every day. In principle it is a large difference, because the boundary of trust
moves from the small kernel to a long chain of compiler, linker, and processor. A certificate built
on top of \texttt{native\_decide} is therefore weaker than a certificate built on top of
\texttt{decide}, even though both print ``accepted'' without complaint.

In the induction theorem above, the list of axioms is not empty either. The command
\texttt{\#print axioms} for \texttt{oddSum\_sq} prints \texttt{[propext, Quot.sound]}. Both are
standard axioms used for classical reasoning and for quotients, and both appear because
\texttt{omega} proves the goal by contradiction inside the integers. This is not an error; it is a
fact that ought to be mentioned when a certificate is used as the basis of a claim.

\section{Limits: rejecting wrong proofs, \texttt{sorry}, and axioms}

\tingkat{advanced}

The ability to reject is as important as the ability to accept. The following file deliberately
contains errors of different kinds, and the output is Lean's answer to each of them.

\begin{verbatim}
-- File: Salah.lean
theorem salah_arit : 2 + 2 = 5 := by decide

theorem salah_rfl (n : Nat) : n + 1 = 1 + n := by rfl

axiom magic : 1 = 2
theorem from_magic : 1 = 2 := magic
#print axioms from_magic

theorem belum_selesai (n : Nat) : n + 0 = n := by sorry
\end{verbatim}

The actual output, cut at the long part:

\begin{verbatim}
Salah.lean:2:37: error: tactic 'decide' proved that the proposition
  2 + 2 = 5
is false
Salah.lean:4:50: error: The rfl tactic failed. Possible reasons:
- The goal is not a reflexive relation ...
- The arguments of the relation are not equal.
...
n : Nat
|- n + 1 = 1 + n
'from_magic' depends on axioms: [magic]
Salah.lean:10:8: warning: declaration uses 'sorry'
\end{verbatim}

Three kinds of rejection appear here, and all three differ in nature. The first is the strongest
rejection: \texttt{decide} does not merely fail to close the goal, it proves that the proposition is
false. For a finite statement, failure and refutation are the same thing. The second is a subtler
rejection: \texttt{rfl} refuses because the two sides are not equal by definition, even though the
statement is true. Here Lean does not blame the statement, it blames the tool in use; the message
even suggests an explicit reflexivity lemma. The third is not a rejection at all.

The line \texttt{axiom magic : 1 = 2} is accepted without complaint. Lean does not check whether the
content of an axiom is reasonable; it only records it.\\index{axiom} After that,
\texttt{from\_magic} is accepted as well, and the whole of arithmetic can be taken apart with
\texttt{False.elim} from there. The only danger sign is the \texttt{\#print axioms} line, which
mentions \texttt{[magic]} calmly, without a warning and without the color red. This is the practical
reason why inspecting the list of axioms is a compulsory habit: the system will not reproach us.

\kamus{axiom}{aksioma}{A rule that we ask to be accepted without proof, like an opening assumption
in a debate; legitimate, but it has to be named.}

The last line of that file uses \texttt{sorry}, and its output is only the warning
\texttt{declaration uses 'sorry'}.\\index{sorry} The file is still processed, the theorem still
``exists'', and anyone who reads the output carefully will see the word \texttt{warning}, not
\texttt{error}. In a large library, one \texttt{sorry} left behind among thousands of theorems is
almost invisible. Lean provides the answer, and the answer is the same as the answer for axioms:
\texttt{\#print axioms} prints \texttt{[sorryAx]}. Lean's mathematics library, namely Mathlib, uses
a stricter policy: a file that contains \texttt{sorry} is rejected by the automatic check before it
is accepted. That policy deserves to be copied by anyone who publishes certificates.

\kamus{sorry}{sorry}{A note saying ``I will do it later'' that makes the work look finished; the
machine accepts it, but records it as a debt.}

Now the most important limit, and the one least often mentioned. Formal verification guarantees that
the term fits the statement. It does not guarantee that the statement fits the intention. If the
specification is too weak, the proof can be perfect and useless.\\index{specification} The following
example shows this in its most naked form: the first statement says ``there is a small perturbation
that does not change the output'', the second statement says ``every small perturbation does not
change the output''. The first can be proved for \emph{all} functions, including the worst function
imaginable, because the witness $\delta=0$ is always available.

\begin{verbatim}
def identitas (n : Nat) : Nat := n

-- weak specification: THERE IS a delta that does not change the output (for all f)
theorem ada_delta (f : Nat -> Nat) (x : Nat) :
    exists delta : Nat, f (x + delta) = f x := by
  exact Exists.intro 0 rfl

-- the strong version is false for f = identity, so it cannot be proved
theorem ada_penyangkal :
    exists x : Nat, exists delta : Nat,
      Not (identitas (x + delta) = identitas x) := by
  exact Exists.intro 0 (Exists.intro 1 (by decide))
\end{verbatim}

Both theorems pass without additional axioms, and that is exactly the problem. The first theorem says
nothing about $f$; it only says that $\delta=0$ is a perturbation that changes nothing, and that
holds for every function, including a classifier that is wrong one hundred percent of the time. A
certificate of this kind passes every check and protects nothing. The second theorem shows the other
side: the statement we actually want, namely the universally quantified version, is false for the
identity function, so no proof can possibly be found.

\begin{tantangan}
A certificate contains one line that looks impossible to be wrong: for every natural number $a$ and
$b$ we have $(a-b)+b=a$. On paper the reason is simple: subtract, add again, and you are back where
you started. Write that line in Lean and close it with \texttt{omega}; the checker rejects it, and
the rejection is not because the tactic is not clever enough.

Two questions. Why is that statement not true for all $a$ and $b$? And what additional condition
makes it true while keeping it checkable by machine?

Hint 1: do not trust the parentheses; ask what happens when $b$ is larger than $a$. Hint 2:
subtraction on the natural numbers stops at zero, so $0-1=0$, hence the left-hand side becomes
$0-1+1=1$ whereas the right-hand side becomes $0$.

If you get stuck, go back to the section ``Tactics: turning the goal into a term'', where $n+1$ and
$1+n$ stop at different forms.
\end{tantangan}

Before reading further, try to answer on your own. The output of this test machine, as it stands, is
the comparison for your answer: the second file contains the same statement with the hypothesis
$b\le a$.

\begin{verbatim}
$ lean t.lean
t.lean:1:49: error: omega could not prove the goal:
a possible counterexample may satisfy the constraints
  d >= 0
  c >= 0
  c - d >= 1
where
 c := b
 d := a
$ lean t2.lean
'potong_benar' depends on axioms: [propext, Quot.sound]
\end{verbatim}

The first output mentions the constraint $b-a\ge 1$, that is a counterexample with $b>a$; at $a=0$
and $b=1$ the left-hand side has the value $1$, the right-hand side $0$. The hypothesis $b\le a$
makes the same line pass, with the two standard axioms that we already know from \texttt{omega}. The
shape of the repair is not an additional tactic, but a condition.

\kertas{The claim that may be made from this chapter: the Lean code fragments that are shown have
been run with Lean 4.10.0 without Mathlib on the machine of writing, and their output is copied as
it stands; Unicode output is reprinted in ASCII and cuts are marked with three dots. What was not
run: the Mathlib tactics (\texttt{linarith}, \texttt{ring}, \texttt{norm\_num},
\texttt{induction'}) and the Mathlib library itself, because it is not installed. What is not proved
here and is not claimed: any certificate about neural networks, language models, or robustness to
attack. Chapter 28 does not contain a single proof of that kind.}

These limits have a direct connection with Chapter 27 and Chapter 29. An honest evaluation of a model
demands that we state the distribution in which the numbers are measured; formal verification
demands that we state the proposition that is proved. Both are the same disease in a different form:
a claim that does not state its conditions. The certificate ``this model is proved safe'' without a
formal statement behind it is as useless as ``99 percent accuracy'' without a note on the test data.
In Chapter 29 and Chapter 30, when we discuss attacks and defenses, the first question that must be
put to every robustness certificate is: how many quantifiers does it have, and in what space do those
quantifiers run?

\section{Worked example}

\tingkat{intermediate}

The numbers in this part are computed, not estimated. The first column is $n$, the second column is
the $n$th odd number that is added at that step, the third column is $S(n)$, and the fourth column is
$n^{2}$.

\begin{center}
\begin{tabular}{rrrr}
\hline
$n$ & term $n$ & $S(n)$ & $n^{2}$ \\
\hline
0  & --  & 0     & 0 \\
1  & 1   & 1     & 1 \\
2  & 3   & 4     & 4 \\
5  & 9   & 25    & 25 \\
10 & 19  & 100   & 100 \\
20 & 39  & 400   & 400 \\
100 & 199 & 10000 & 10000 \\
\hline
\end{tabular}
\end{center}

The value $S(100)=10000$ is worth checking by hand once, because it shows that the formula is not a
coincidence for small numbers: the sum of a hundred odd numbers, from $1$ to $199$, equals
$100\cdot 100$. A check with Python over all values $n<5000$ gives \texttt{True}, and a check of the
algebraic identity $(k+1)^{2}=k^{2}+2k+1$ for all $k<100000$ also gives \texttt{True}. Neither check
is a proof, but both rule out the possibility that the formula being proved has a typo.

There is one more number that is more informative than the table above, namely the size of the proof
after elaboration. The tactic script for \texttt{oddSum\_sq} is eight lines long. When Lean prints
the term that results from elaborating it with \texttt{\#print}, the result is 10,783 characters in
149 lines. That means one line of tactics expands on average into about nineteen lines of term. No
human reads that term; what reads it is the kernel, and that is the division of labor that is the
heart of the whole method.

The last number concerns the limit of direct checking. In the default setting, the recursion-depth
limit has the value 1000. With that limit, the tactic \texttt{decide} succeeds in proving
$S(n)=n^{2}$ only for $n\le 122$. The first experiment that fails is $n=123$, and the message that
appears mentions that the maximum recursion depth has been reached. If the limit is raised to 100000,
the value $n=2000$ is proved successfully in 0.48 seconds. At $n=20000$ the check fails because the
recursion is too deep for the stack space. So direct checking is not an alternative that scales,
whereas an inductive proof does not care about the size of $n$.

\section{Computational lab}

\tingkat{advanced}

This lab runs two things: the Lean file that is the same as the one discussed above, and a
cross-check with Python. The Lean file is written without \texttt{import}, so it uses only the
kernel.

\begin{lab}
\begin{verbatim}
-- BilanganGanjil.lean : Lean 4.10.0, without Mathlib
def oddSum : Nat -> Nat
  | 0     => 0
  | n + 1 => oddSum n + (2 * n + 1)

theorem oddSum_sq (n : Nat) : oddSum n = n * n := by
  induction n with
  | zero => rfl
  | succ k ih =>
    rw [oddSum]
    rw [ih]
    simp [Nat.succ_eq_add_one, Nat.mul_succ, Nat.add_mul, Nat.one_mul]
    omega

#eval oddSum 0
#eval oddSum 5
#eval oddSum 10
#eval oddSum 100
#print axioms oddSum_sq
\end{verbatim}

\begin{verbatim}
def odd_sum(n):
    return sum(2 * i + 1 for i in range(n))

print([odd_sum(n) for n in range(11)])
print(odd_sum(20), odd_sum(100))
print(all(odd_sum(n) == n * n for n in range(5000)))
print(all((k + 1) * (k + 1) == k * k + 2 * k + 1 for k in range(100000)))
\end{verbatim}
\end{lab}

The commands and the output, as they stand:

\begin{verbatim}
$ lean --version
Lean (version 4.10.0, x86_64-unknown-linux-gnu, commit c375e19f6b65, Release)
$ time lean BilanganGanjil.lean
0
25
100
10000
'oddSum_sq' depends on axioms: [propext, Quot.sound]
time 0.14 s
$ python3 BilanganGanjil.py
[0, 1, 4, 9, 16, 25, 36, 49, 64, 81, 100]
400 10000
True
True
\end{verbatim}

The first four lines of the Lean output are the results of \texttt{\#eval}, and all four match the
table in the previous part. The fifth line is the list of axioms, and its contents must be read
carefully: \texttt{propext} and \texttt{Quot.sound} are not an error, but a consequence of the way
\texttt{omega} works. The Python output confirms the same thing from the other direction: the list of
the first eleven values is a list of perfect squares, $S(100)=10000$, and both exhaustive checks
print \texttt{True}. The time 0.14 seconds is the time of this test machine, and a number of that
kind depends on the hardware; what is meaningful is not the number, but the fact that the check
finishes before we have time to read the file.

A reader who wants to repeat this lab needs two things. First, Lean 4, which can be installed through
\texttt{elan}, the official version manager of the Lean project; after installation, the command
\texttt{lean --version} must print a version. Second, the file above is saved as
\texttt{BilanganGanjil.lean} and run with \texttt{lean BilanganGanjil.lean} from the same directory;
no \texttt{lake} project is needed because there is no import. To use \texttt{linarith} or
\texttt{ring}, the reader needs a project with Mathlib, and from that moment \texttt{lean} alone is
not enough; the command becomes \texttt{lake env lean}. This difference is often the first source of
confusion for a beginner, and mentioning it is better than letting it be discovered through an
\texttt{unknown tactic} message.

\begin{resep}
\begin{enumerate}
  \item Create a plain text file named \texttt{Uji.lean}; that extension is what makes the command
        \texttt{lean} recognize it as a source file.
  \item Write all \texttt{import}s on the topmost lines, one per line, before any declaration.
  \item Write the statement first, complete with the type of each variable, then the proof after
        \texttt{:= by}.
  \item Run \texttt{lean Uji.lean} from the directory of that file and read the output from the first
        line; the first error is usually the real one.
  \item Ask for \texttt{\#print axioms name\_of\_theorem} before calling the result proved.
\end{enumerate}
A warning about order: \texttt{import} must not be preceded by a declaration. One \texttt{def}
slipped above it makes Lean reject the whole file with the message
\texttt{invalid 'import' command}, and the file is not checked at all. The correct order is not the
author's taste; it is a condition for the file to be checkable.
\end{resep}

\begin{untukdosen}
The teaching order that works in class: start from the failure of \texttt{rfl} on $n+1=1+n$, not
from the definition. Students who see that something ``obvious'' is rejected will ask, and that
question is the doorway into Curry--Howard. Only after that come the definition of a proposition as a
type, then the first three tactics, then induction. Leave one full session for \texttt{\#print
axioms}; that is where students change from users into checkers.

Two misconceptions appear most often. First, ``Lean proves that my program is correct''. What is
true: Lean checks that a term has a certain type of statement. If the statement does not mention the
important thing, for example a quantifier over perturbations, the certificate is empty. Second,
``\texttt{simp} is a proof''. \texttt{simp} is a procedure whose output depends on the active
collection of lemmas, so a proof that depends on it can stop working when the library changes; that
is why we train students to write the list of lemmas explicitly.

A blackboard activity that works: ask two groups to write a specification for ``a classifier that is
robust to perturbation''. One group will almost always write an existential quantifier. Put both
formulations on the board, then prove the weak one for a constant function. The class will see for
itself that the sentence that is proved and the sentence that is meant are two different sentences.

A sign of readiness to go on: students can explain why $n+1=1+n$ requires a lemma, can write an
inductive proof without copying, and can name one axiom that their own theorem uses.
\end{untukdosen}

\section{Summary}

\begin{ringkasan}
\begin{itemize}
  \item An informal proof convinces a trained reader; a formal proof is a term whose type is its
        statement, checked by a small kernel that knows no mathematics.
  \item According to Curry--Howard, implication is a function type, conjunction is a pair, the
        universal quantifier is a dependent function, and $\bot$ is uninhabited.
  \item Tactics change the goal state; \texttt{intro}, \texttt{exact}, \texttt{apply},
        \texttt{rw}, \texttt{cases}, and \texttt{induction} are enough for basic arithmetic.
  \item Induction proves a statement for all $n$ with two goals, not by checking all $n$;
        \texttt{decide} by contrast stops scaling very quickly.
  \item \texttt{\#print axioms} is the main instrument of inspection. The tactic
        \texttt{native\_decide} leaves behind the axiom \texttt{Lean.ofReduceBool}. The tactic
        \texttt{sorry} leaves behind \texttt{sorryAx}. A homemade axiom leaves behind its own name.
  \item Verification guarantees that the term fits the statement, not that the statement fits the
        intention; a weak specification produces an empty certificate.
\end{itemize}
\end{ringkasan}

\section{Exercises}

\begin{latihan}
\begin{enumerate}
  \item Assess this claim: \emph{if a Lean file is accepted without an error message, then all the
        statements inside it are true}. Decide true, false, or not enough evidence, and name one
        command that can be used to check that claim for a particular file.

  \item Prove in Lean, without Mathlib, that for every $n$ we have $2\cdot S(n)=n\cdot(2n)$ with $S$
        as in this chapter. State first the auxiliary lemma that you need.

  \item Prove that the sum of the first $n$ even numbers, namely $0+2+\dots+(2n-2)$, equals
        $n(n-1)$. Write the informal proof, then write the Lean version and name the tactic that is
        used at every step.

  \item Compute with Python: the sum of the first 37 odd numbers, the sum of the first 250 odd
        numbers, and check that both are perfect squares. Include the code and its output.

  \item Compute with Python how many cases must be checked if one wants to prove $S(n)=n^{2}$ by
        direct checking for all $n\le 1000$, then compare it with the number of goals that an
        inductive proof generates.

  \item Explain why \texttt{rfl} closes the goal $n+0=n$ but does not close $0+n=n$. The hint lies
        in the way addition on the natural numbers is defined recursively; name which argument is
        reduced first.

  \item Write a weak specification and a strong specification for the claim ``this classifier does
        not change under a perturbation of size at most one''. Prove the weak version in Lean for an
        arbitrary function, then explain why the strong version cannot be proved for the identity
        function.

  \item Inspect with \texttt{\#print axioms} three theorems that you wrote yourself: one using
        \texttt{decide}, one using \texttt{omega}, one using \texttt{sorry}. Report the three lists
        of axioms and explain the origin of every name.

  \item A certificate contains \texttt{theorem aman : exists delta : Nat, f (x + delta) = f x} with
        \texttt{delta} as a free variable in the context. Explain two errors that may be present in
        that statement, and write the version that you consider correct.

  \item Prove in Lean that $\neg(P\wedge\neg P)$ for every proposition $P$. Hint: start with
        \texttt{intro h}, then take the components with \texttt{h.1} and \texttt{h.2}, and use the
        second component on the first one.
\end{enumerate}

Hints for the three hardest items. Item 2: the auxiliary lemma that is needed has the shape
$2\,(k\cdot k+(2k+1))=(k+1)\cdot(2(k+1))$, and after \texttt{rw} on the definition and on the
induction hypothesis, the remaining goal is linear so that \texttt{omega} is enough. Item 7: in the
weak version, the witness $\delta=0$ makes both sides equal by definition, so \texttt{rfl} closes it
immediately; in the strong version, the identity function gives a counterexample with $\delta=1$.
Item 9: note that the free \texttt{delta} can be chosen by the reader of the statement, not by the
attacker, so that statement asserts no robustness at all; the second error is the absence of a
quantifier over $x$.
\end{latihan}

\begin{untukpeneliti}
What is really proved in this chapter is only small arithmetic statements: $S(n)=n^{2}$ for all
natural numbers, a few propositional tautologies, and two statements about specifications that are
too weak. All of them are checked by the Lean 4.10.0 kernel without Mathlib, and their lists of
axioms are reported. What is popularized far beyond that: the phrase ``formally verified'' is often
used for ``checked by an automatic tool'' or even for ``tested exhaustively''. Only the first
produces an object that a third party can recheck.

Primary sources. Avigad, de Moura, Kong, and Ullrich, \emph{Theorem Proving in Lean 4}, the sections
``Propositions and Proofs'' (propositions as types, proof irrelevance, the axiom warning),
``Tactics'' (the goal state), and ``Induction and Recursion'' (\texttt{induction},
\texttt{sorryAx}). Avigad and Massot, \emph{Mathematics in Lean}, Chapters 2, 3.1, 5.2, and 8; the
Indonesian edition uses the same section numbering.

Open questions. This chapter does not contain a formal proof about neural networks. An interesting
robustness specification has the shape
$\forall x\,\forall\delta,\ \lVert\delta\rVert\le\varepsilon\to f(x+\delta)=f(x)$, that is a
statement over a continuous space with a function with hundreds of thousands of parameters.
Formalizing it demands real arithmetic that is decidable at that scale, and the Lean libraries today
do not provide it for trained models. The claim ``we verified the robustness of the model'' in a
paper usually refers to verification with interval abstraction or with SMT, not to a Lean proof; both
are legitimate, but their guarantees differ and must not be called by the same word.

Limits of validity. All the claims in this chapter hold for Lean 4.10.0 with its default kernel, and
the details of the output can change between versions. What does not change is the structure of the
argument: statement, term, kernel, list of axioms.
\end{untukpeneliti}

\peta{This chapter follows Avigad, de Moura, Kong, and Ullrich, \emph{Theorem Proving in Lean 4}, the
sections ``Propositions and Proofs'' (propositions as types, Curry--Howard, proof irrelevance, the
axiom warning), ``Quantifiers and Equality'', ``Tactics'' (the goal state, \texttt{intro},
\texttt{apply}, \texttt{exact}), and ``Induction and Recursion'' (\texttt{induction}, \texttt{sorry}
and \texttt{sorryAx}). The counterpart in Avigad and Massot, \emph{Mathematics in Lean}, is in
Chapter 2 (Basics, in particular 2.3 Using Theorems and Lemmas and 2.4 More Examples Using apply and
rw), Chapter 3.1 (Implications and Universal Quantifiers), Chapter 5.2 (Induction and Recursion), and
Chapter 8 (Hierarchies) for typeclasses; the Indonesian edition keeps the same section numbering. For
the connection with evaluation practice, see Chapter 27 of this book; for attacks that make
robustness specifications nontrivial, see Chapter 27 and Chapter 28.}
